
\chapter{Jean-Fran\c{c}ois Le Gall (2019)}
\index{Le Gall, Jean-Fran\c{c}ois}

\begin{quote}
\it ``Jean-Fran\c cois Le Gall has made transformative contributions to probability theory. His work on the Brownian snake, random trees, and random maps has revealed profound connections between stochastic processes and geometry. The construction of the Brownian map is one of the most remarkable achievements in modern probability.'' --- Wolf Prize Citation, 2019
\end{quote}

\section{Early Life and Career}

Jean-Fran\c{c}ois Le Gall was born in 1959 in Paris, France. He studied at the University of Paris-Sud (Orsay), receiving his Ph.D.\ in 1982 under the supervision of Claude Dellacherie. He has spent his career at the University of Paris-Sud (now Universit\'{e} Paris-Saclay) and the \'Ecole Normale Sup\'{e}rieure, where he is currently a professor. He was awarded the Wolf Prize in Mathematics in 2019.

\section{Main Contributions}

Jean-Fran\c{c}ois Le Gall (born 1959) is a French mathematician whose work in probability theory has opened new connections between stochastic processes, geometry, and combinatorics. He was awarded the Wolf Prize in Mathematics in 2019 alongside Gregory Lawler.

The Wolf Prize citation reads: ``for their groundbreaking contributions to probability theory and mathematical physics.''

Le Gall's core contributions include:
\begin{itemize}
    \item \textbf{The Brownian Map:} The construction and proof that the scaling limit of random planar maps (random surfaces discretized by polygons) is a universal continuum random surface called the Brownian map, which is homeomorphic to the 2-sphere.
    \item \textbf{The Brownian Snake:} A path-valued stochastic process that describes the genealogy of the continuum random tree and encodes the distances in the Brownian map.
    \item \textbf{The Continuum Random Tree (CRT):} Foundational work on the CRT as a universal scaling limit for random trees, now a fundamental object in probability.
    \item \textbf{Superprocesses:} Deep results on measure-valued branching processes (Dawson--Watanabe superprocesses), including connections to partial differential equations.
    \item \textbf{Stochastic Calculus:} The Le Gall--Yor theorem on the representation of Bessel processes and the study of stochastic differential equations.
    \item \textbf{Random Fractals:} The computation of the Hausdorff dimension of the trace of the Brownian snake and related random sets.
\end{itemize}

\section{Origin of the Problem}

\subsection{Random Planar Maps}

A planar map is a graph embedded in the sphere. Random planar maps, sampled uniformly from all maps of a given size, were studied in theoretical physics as models of random surfaces (2D quantum gravity).

The fundamental question: what does a large random planar map look like at large scales? Does it have a continuum limit (scaling limit) analogous to the convergence of simple random walk to Brownian motion?

\subsection{The Brownian Snake and Genealogies}

The Brownian snake is a stochastic process whose state is a finite path (a continuous function from $[0,t]$ to $\R$). It was introduced to describe the genealogy of the branching process known as the ``super-Brownian motion.'' The snake encodes the entire family tree: the paths represent the spatial trajectories of individuals in the population.

\section{How It Was Solved and the Intuitions/Tools Needed}

\subsection{The Construction of the Brownian Map}

The Brownian map is constructed from two independent objects:
\begin{enumerate}
    \item A continuum random tree (the CRT) encoded by a normalized Brownian excursion.
    \item A random ``label'' or ``height'' function on the tree, given by the Brownian snake.
\end{enumerate}

The key insight: the Brownian map is a quotient of the CRT by an equivalence relation that identifies points with the same label and the same distance structure. Le Gall proved that this metric space is homeomorphic to the 2-sphere and has Hausdorff dimension 4.

\subsection{Tools and Techniques}

\begin{itemize}
    \item \textbf{Excursion Theory:} The Brownian excursion describes the contour of the CRT.
    \item \textbf{The Brownian Snake:} Labels are assigned to tree points using the snake process.
    \item \textbf{Gromov--Hausdorff Convergence:} The space of compact metric spaces with the Gromov--Hausdorff topology provides the framework for convergence.
    \item \textbf{Combinatorial Enumeration:} Exact formulas for the number of planar maps (Tutte's formulas) provide the discrete input.
    \item \textbf{Scaling Limits:} Convergence in distribution of rescaled discrete objects to continuous limits.
\end{itemize}

\section{Mathematical Theory}

\subsection{The Continuum Random Tree}

\begin{definition}[Continuum Random Tree]
The \textbf{continuum random tree} (CRT) is a random compact metric tree encoded by a normalized Brownian excursion $(e_t)_{0 \leq t \leq 1}$:
\[
d_e(s,t) = e_s + e_t - 2 \inf_{u \in [s,t]} e_u
\]
The tree is $\mathcal{T}_e = [0,1]/\sim$ where $s \sim t$ if $d_e(s,t) = 0$.
\end{definition}

\begin{theorem}[Aldous, 1991; Le Gall, 1993]
The CRT is the universal scaling limit of many families of random trees, including uniform random plane trees, random binary trees, and Galton--Watson trees conditioned on total size.
\end{theorem}

\subsection{The Brownian Snake}

\begin{definition}[Brownian Snake]
The \textbf{Brownian snake} is a continuous stochastic process $(W_s)_{s \geq 0}$ where each $W_s$ is a finite path (a continuous function $[0,\zeta_s] \to \R$). The process is characterized by:
\begin{enumerate}
    \item The lifetime process $(\zeta_s)_{s \geq 0}$ is a reflected Brownian motion.
    \item Conditionally on $\zeta$, the spatial motion is a Markov process: when $\zeta$ increases, the path is extended by a Brownian motion; when $\zeta$ decreases, the path is erased from the tip.
\end{enumerate}
\end{definition}

\subsection{The Brownian Map}

\begin{definition}[Brownian Map]
The \textbf{Brownian map} is the random metric space $(\mathcal{S}_L, D_L)$ obtained by:
\begin{enumerate}
    \item Taking the CRT $\mathcal{T}_e$ with labels $(Z_v)_{v \in \mathcal{T}_e}$ from the Brownian snake.
    \item Defining the pseudometric:
    \[
    D^*(x,y) = \inf_{x = v_0, v_1, \ldots, v_k = y} \sum_{i=0}^{k-1} d(x_i, x_{i+1})^{\natural}
    \]
    where $d(x,y)^{\natural} = Z_x + Z_y - 2\max_{w \in [x,y]} Z_w$.
    \item Taking the quotient $\mathcal{S}_L = \mathcal{T}_e / \sim$ where $x \sim y$ if $D^*(x,y) = 0$.
\end{enumerate}
\end{definition}

\begin{theorem}[Le Gall, 2007; Le Gall--Paul, 2008]
Let $M_n$ be a uniformly random quadrangulation of the sphere with $n$ faces, viewed as a metric space with distances rescaled by $n^{-1/4}$. Then as $n \to \infty$:
\[
M_n \xrightarrow{d} \mathcal{S}_L
\]
in the Gromov--Hausdorff sense, where $\mathcal{S}_L$ is the Brownian map. Moreover, $\mathcal{S}_L$ is almost surely homeomorphic to the 2-sphere and has Hausdorff dimension 4.
\end{theorem}

\subsection{Super-Brownian Motion}

\begin{definition}[Super-Brownian Motion]
\textbf{Super-Brownian motion} $(X_t)_{t \geq 0}$ is a measure-valued Markov process where $X_t$ is a random finite measure on $\R^d$. It satisfies the branching property: the total mass $M_t = X_t(\R^d)$ is a continuous-state branching process.
\end{definition}

Le Gall established the deep connection between super-Brownian motion and the nonlinear PDE:
\[
\partial_t u = \frac{1}{2}\Delta u + \beta u^2
\]
through the Brownian snake representation.

\subsection{The Feynman--Kac Formula}

The Feynman--Kac formula provides the fundamental link between Brownian motion and parabolic partial differential equations. It expresses solutions of the heat equation with a potential as expectations of Brownian motion functionals.

\begin{theorem}[Feynman--Kac]
Let $V \colon \R^d \to \R$ be a bounded measurable function and let $(B_t)_{t \geq 0}$ be a standard Brownian motion in $\R^d$ started at $x$. Define
\[
    u(t,x) = \E_x\!\left[\exp\!\left(-\int_0^t V(B_s)\,ds\right) f(B_t)\right]
\]
for a bounded measurable function $f$. Then $u$ is the unique bounded measurable solution of the Cauchy problem
\begin{align}
    \partial_t u &= \tfrac{1}{2}\Delta u - V\, u, \qquad t > 0,\; x \in \R^d, \label{eq:feynman-kac-pde}\\
    u(0,x) &= f(x). \label{eq:feynman-kac-ic}
\end{align}
\end{theorem}

\begin{remark}
The Feynman--Kac formula can be reformulated in terms of the semigroup $(T_t)_{t \geq 0}$ associated to the operator $\mathcal{L} = \frac{1}{2}\Delta - V$. For each $t \geq 0$, the operator $T_t$ acts on $L^2(\R^d)$ by
\[
    (T_t \varphi)(x) = \E_x\!\left[\exp\!\left(-\int_0^t V(B_s)\,ds\right) \varphi(B_t)\right].
\]
The family $(T_t)_{t \geq 0}$ forms a strongly continuous contraction semigroup on $L^2(\R^d)$, and its infinitesimal generator is $\mathcal{L}$ with domain $H^2(\R^d)$.
\end{remark}

\begin{remark}[Heat Kernel and Transition Density]
When $V \equiv 0$, the semigroup reduces to the free heat semigroup, and the solution is given explicitly by the heat kernel:
\[
    (T_t \varphi)(x) = \int_{\R^d} p_t(x,y)\,\varphi(y)\,dy,
\]
where
\[
    p_t(x,y) = \frac{1}{(2\pi t)^{d/2}}\,\exp\!\left(-\frac{|x-y|^2}{2t}\right)
\]
is the Gaussian transition density of Brownian motion. More generally, if $V$ is a bounded potential, the Feynman--Kac semigroup admits a smooth kernel $p_t^V(x,y)$ satisfying
\[
    \partial_t p_t^V = \tfrac{1}{2}\Delta\, p_t^V - V\, p_t^V
\]
with $p_0^V(x,y) = \delta_y(x)$. This kernel satisfies the Chapman--Kolmogorov equation $p_{s+t}^V(x,y) = \int p_s^V(x,z)\,p_t^V(z,y)\,dz$ and is strictly positive for all $t > 0$.
\end{remark}

\subsection{Brownian Motion as a Scaling Limit of Random Walks}

Brownian motion arises canonically as the scaling limit of simple random walks. This invariance principle, due to Donsker, is one of the cornerstones of modern probability.

\begin{theorem}[Donsker's Invariance Principle]
Let $(S_n)_{n \geq 0}$ be a simple symmetric random walk on $\Z^d$, so $S_n = \xi_1 + \cdots + \xi_n$ where the $\xi_i$ are i.i.d.\ with $\P(\xi_i = e_j) = \P(\xi_i = -e_j) = \frac{1}{2d}$ for each coordinate direction $e_j$. Define the continuous-time interpolation $W^{(n)}$ on $[0,1]$ by
\[
    W^{(n)}_t = \frac{1}{\sqrt{n}}\,S_{\lfloor nt \rfloor} + \frac{nt - \lfloor nt \rfloor}{\sqrt{n}}\,\xi_{\lfloor nt \rfloor + 1}.
\]
Then $W^{(n)}$ converges in distribution in $C([0,1],\R^d)$ to a standard Brownian motion $(B_t)_{0 \leq t \leq 1}$ equipped with the uniform topology.
\end{theorem}

\begin{remark}
The functional central limit theorem (FCLT) extends Donsker's principle to random walks with increments having mean zero and finite variance $\sigma^2$. In this case, the limit is $\sigma$ times a standard Brownian motion. The convergence holds in the Skorokhod space $D([0,1],\R^d)$ for general increments and in $C([0,1],\R^d)$ when the walk is continuous (lattice-valued).
\end{remark}

\begin{remark}[Scaling Exponents]
The scaling relation $W^{(n)}_t \approx n^{-1/2} S_{nt}$ has important generalizations. For random planar maps, the analogous scaling involves $n^{-1/4}$ in the metric, reflecting the fact that the Brownian map has Hausdorff dimension 4 while Brownian motion has dimension 2. This dimensional doubling is a hallmark of the connection between random walks on planar maps and the Brownian snake: the genealogical tree of the walk (a CRT) has dimension 2, but the two-dimensional surface structure encoded by the Brownian snake raises the dimension to 4.
\end{remark}

\subsection{The Brownian Plane}

The Brownian plane, introduced by Le Gall and Miermont (2011), is the scaling limit of uniformly random planar maps viewed as random metric spaces. It plays the role that Brownian motion plays for random walks: the universal random metric space underlying two-dimensional random geometry.

\begin{definition}[Brownian Plane]
The \textbf{Brownian plane} is the random metric space $(\mathcal{Q}, d_{\mathcal{Q}})$ defined as follows. Let $(e_t)_{0 \leq t \leq 1}$ be a normalized Brownian excursion, and let $(Z_v)_{v \in \mathcal{T}_e}$ be the labels of the Brownian snake indexed by the CRT $\mathcal{T}_e$. For $x, y \in \mathcal{T}_e$, define the pseudodistance
\[
    d^*(x,y) = Z_x + Z_y - 2\inf_{w \in [x,y]_{\mathcal{T}_e}} Z_w,
\]
where $[x,y]_{\mathcal{T}_e}$ is the geodesic path in the CRT. The Brownian plane is obtained by the Poisson construction: embed the CRT into the real line via the contour function, then assign a mass measure to each edge of the tree proportional to the local time of the Brownian motion. The metric $d_{\mathcal{Q}}$ is then defined as the quotient metric on $\mathcal{T}_e / \!\sim$ where $x \sim y$ if $d^*(x,y) = 0$.
\end{definition}

\begin{theorem}[Le Gall--Miermont, 2011]
Let $M_n$ be a uniformly random planar map (e.g., a triangulation or quadrangulation of the sphere) with $n$ edges, viewed as a metric space with the graph distance rescaled by $n^{-1/2}$. Then as $n \to \infty$:
\[
    (M_n, n^{-1/2} d_{\text{graph}}) \xrightarrow{d} (\mathcal{Q}, d_{\mathcal{Q}})
\]
in the Gromov--Hausdorff sense. The Brownian plane is an infinite random metric space with the following properties:
\begin{enumerate}
    \item It is almost surely a geodesic metric space.
    \item It is isometry-invariant: the law of $(\mathcal{Q}, d_{\mathcal{Q}})$ is invariant under the action of isometries fixing a marked point.
    \item The Hausdorff dimension of $\mathcal{Q}$ is almost surely 4.
    \item The spectral dimension is $\frac{4}{3}$.
\end{enumerate}
\end{theorem}

\begin{remark}
The Brownian plane is closely related to the Brownian map but differs in an important way: the Brownian plane is infinite (non-compact), while the Brownian map is compact. The Brownian map can be obtained from the Brownian plane by a ``sphereification'' procedure that identifies the boundary at infinity. Formally, the Brownian map is the quotient of a suitably defined compactification of the Brownian plane.
\end{remark}

\begin{remark}[Coordinates]
The Brownian plane can be described through its coordinate process: a real-valued process $(B_t)_{t \geq 0}$ coupled with a tree-valued process (the CRT) and a label process from the Brownian snake. The geodesic distance between two points with contour times $s, t \in [0,1]$ is
\[
    d_{\mathcal{Q}}(s,t) = \sup_{r \in [s,t]} \left( Z_r - \inf_{u \in [s \wedge r, t \vee r]} Z_u \right),
\]
where $Z$ is the Brownian snake label process. This self-similar structure underpins the universality of the Brownian plane as a scaling limit.
\end{remark}

\subsection{Applications to Random Geometry}

Le Gall's work on the Brownian map, the Brownian plane, and related processes has had profound applications in random geometry and combinatorial probability.

\subsubsection{Random Trees and the CRT}

The continuum random tree serves as the universal scaling limit of a broad family of discrete random trees. Le Gall's results on the CRT as a scaling limit of Galton--Watson trees conditioned on their total population size have become foundational.

\begin{proposition}[Le Gall, 1993]
Let $(T_n)$ be a Galton--Watson branching process with offspring distribution having finite variance $\sigma^2 > 0$, conditioned on $|T_n| = k$. Then the rescaled tree $(k^{-1/2} d_{T_n})$ converges in distribution in the Gromov--Hausdorff topology to $\sigma \cdot \mathcal{T}_e$, where $\mathcal{T}_e$ is the CRT encoded by a normalized Brownian excursion.
\end{proposition}

The CRT is the ``tree of geodesics'' of the Brownian map: the map is constructed by attaching labels to the CRT via the Brownian snake and then quotienting. Random trees arising in computer science (binary search trees, random recursive trees) and in physics (spanning trees, minimum spanning trees) all converge to the CRT after appropriate rescaling, making it a universal object.

\subsubsection{The Mated-Cauchy Process and L\'evy Trees}

Le Gall's theory extends beyond Brownian motion to processes with jumps. The mated-Cauchy process $(X_t)_{t \geq 0}$ is a symmetric $\alpha$-stable L\'evy process in $\R^d$ (with $\alpha = 1$, i.e., Cauchy increments) whose genealogical structure is described by a L\'evy tree rather than the CRT. Le Gall and Le Jan showed that the exploration process of the branching structure associated to a superprocess with stable spatial motion is described by a L\'evy snake, a generalization of the Brownian snake.

\begin{remark}
The Hausdorff dimension of the range of the L\'evy snake in the superprocess setting is $\alpha$ (the stability index), and the corresponding random metric space (the ``Brownian map'' analogue for $\alpha$-stable processes) has dimension $2d/\alpha$ when $d < \alpha$. This provides a family of universal random geometries parameterized by $\alpha$, with the Brownian case $\alpha = 2$ recovering dimension 4 for $d = 2$.
\end{remark}

\subsubsection{Random Planar Maps and Quantum Gravity}

The convergence of random planar maps to the Brownian map provides the rigorous mathematical foundation for two-dimensional quantum gravity. In physics, the partition function of 2D quantum gravity is defined as a sum over equivalence classes of metrics on the two-sphere. Random planar maps serve as a discrete regularization of this sum.

Le Gall's construction shows that the continuum limit is the Brownian map, which is homeomorphic to $\mathbb{S}^2$ and has Hausdorff dimension 4. This result confirmed conjectures from the physics literature (originally due to Chekhov, David, Knizhnik, and Polyakov) and established the ``multi-point loop gas'' interpretation of the Brownian map.

\begin{theorem}[Universality of the Brownian Map]
The Brownian map is the scaling limit not only of random quadrangulations but of any ``nice'' family of planar maps, including random triangulations, Eulerian maps, and maps with prescribed face degrees satisfying a moment condition. Specifically, if $(M_n)$ is a family of planar maps with $n$ faces such that
\[
    \E\!\left[\text{(number of faces of degree }k\text{)}^2\right] \leq C\, n
\]
for all $k \geq 3$ and some constant $C$, then $(M_n, n^{-1/4} d_{\text{graph}})$ converges to the Brownian map in the Gromov--Hausdorff topology.
\end{theorem}

\subsubsection{The mated-Cauchy Process and Random Geometry in Higher Dimensions}

Le Gall's framework for the Brownian map has been extended to higher-dimensional settings by several authors, including Abraham, Delmas, and Hoscheit. The key insight is that the Brownian snake construction generalizes to higher-dimensional spatial motions: the ``Brownian map'' in $d$ dimensions has Hausdorff dimension $d+2$ and arises as the scaling limit of random maps with higher-dimensional face structure. Le Gall's original techniques---excursion theory, the Brownian snake, and Gromov--Hausdorff convergence---remain the essential tools.

\begin{remark}
The study of random geometry in dimensions $d \geq 3$ remains one of the major open problems in probability. While the $d = 2$ case (the Brownian map) is completely understood, the analogous objects in $d \geq 3$ are conjectured to exist but their construction requires fundamentally new ideas beyond the Brownian snake framework.
\end{remark}

\section{Impact on Science and Technology}

\subsection{In Mathematics}

\begin{enumerate}
    \item \textbf{Random Geometry:} The Brownian map is the universal continuum limit of random planar maps, providing the mathematical foundation for 2D quantum gravity.
    \item \textbf{Probability Theory:} The Brownian snake and CRT have become standard tools.
    \item \textbf{Combinatorics:} Connections between planar map enumeration and random matrix theory.
    \item \textbf{Stochastic Processes:} The Le Gall--Yor theorem and results on Bessel processes are standard.
\end{enumerate}

\subsection{In Physics}

\begin{enumerate}
    \item \textbf{Quantum Gravity:} Random planar maps are models of Euclidean quantum gravity in two dimensions.
    \item \textbf{String Theory:} The continuum limit of random surfaces connects to string theory.
    \item \textbf{Statistical Mechanics:} Random maps are related to the Ising model, percolation, and other models.
\end{enumerate}

\section{Frequently Asked Questions}

\subsection*{Q1: What is the Brownian map?}
The Brownian map is the universal continuum random surface that arises as the scaling limit of large random planar maps. It is homeomorphic to the 2-sphere and has Hausdorff dimension 4.

\subsection*{Q2: What is the Brownian snake?}
A stochastic process whose state is a finite Brownian path. It encodes the genealogy and spatial motion of a branching population and is the key tool for constructing the Brownian map.

\subsection*{Q3: What is the continuum random tree?}
A random metric tree that arises as the scaling limit of random discrete trees. It is encoded by a Brownian excursion and is the tree structure underlying the Brownian map.

\subsection*{Q4: What are random planar maps used for?}
They are discrete models of random surfaces. In physics, they model 2D quantum gravity. In probability, they are the discrete objects whose scaling limit is the Brownian map.

\subsection*{Q5: What should an undergraduate study to understand Le Gall's work?}
Probability theory (stochastic processes, Brownian motion, martingales), measure theory, and basic topology. Recommended: \textit{Brownian Motion and Stochastic Calculus} (Karatzas--Shreve), \textit{Random Trees} (Evans), and \textit{Brownian Maps} (Le Gall).

\section{Related Chapters}
See also Chapter~17 (It\^o) on stochastic calculus.

\section*{Exercises for the Reader}
\addcontentsline{toc}{section}{Exercises}

\begin{enumerate}
    \item [I] Show that the CRT encoded by a Brownian excursion is a random metric tree.
    \item [A] Compute the Hausdorff dimension of the Brownian snake path.
    \item [I] Show that the number of planar quadrangulations with $n$ faces grows as $c \cdot 12^n n^{-7/2}$.
    \item [A] Prove that the Brownian map is a geodesic metric space.
    \item [A] Compute the spectral dimension of the CRT.
    \item [I] Show that super-Brownian motion satisfies the branching property.
    \item [A] Prove that the Brownian map has Hausdorff dimension 4.
\end{enumerate}

\section*{Timeline}
\addcontentsline{toc}{section}{Timeline}

\begin{tabularx}{\linewidth}{|p{1.5cm}|>{\raggedright\arraybackslash}X|}
\hline
\textbf{Year} & \textbf{Event} \\ \hline
1959 & Born in Morlaix, France \\ \hline
1982 & \'Ecole Normale Sup\'erieure \\ \hline
1987 & Ph.D.\ from Universit\'e Paris VI (Pierre et Marie Curie) \\ \hline
1993 | Professor at Universit\'e Paris-Sud (Orsay) \\ \hline
1999 | Brownian snake and superprocesses monograph \\ \hline
2005 | Moves to Universit\'e Paris-Saclay \\ \hline
2007 | Construction of the Brownian map \\ \hline
2008 | Brownian map is the sphere (with Paul) \\ \hline
2013 | Lo\`eve Prize \\ \hline
2019 | Wolf Prize in Mathematics \\ \hline
\end{tabularx}

\section*{Glossary}
\addcontentsline{toc}{section}{Glossary}

\begin{description}
    \item[Brownian map] The universal continuum random surface, limit of random planar maps.
    \item[Brownian snake] A path-valued process encoding the genealogy of branching processes.
    \item[Continuum random tree] A random metric tree encoded by a Brownian excursion.
    \item[Gromov--Hausdorff convergence] A notion of convergence for metric spaces.
    \item[Planar map] A graph embedded in the sphere, considered up to homeomorphism.
    \item[Quadrangulation] A planar map whose faces are all quadrilaterals.
    \item[Scaling limit] The continuum limit of a rescaled discrete random object.
    \item[Superprocess] A measure-valued branching process.
\end{description}

\section{Citations}\subsection*{Primary Sources}
\begin{enumerate}
    \item J.-F.\ Le Gall, \textit{The Brownian map is the scaling limit of uniform random plane quadrangulations}, Acta Math.\ 210 (2013), 319--401.
    \item J.-F.\ Le Gall, \textit{The topological structure of scaling limits of random planar maps}, Invent.\ Math.\ 169 (2007), 621--670.
    \item J.-F.\ Le Gall, \textit{Spatial Branching Processes, Random Snakes, and Partial Differential Equations}, Lectures in Math.\ ETH Z\"urich, Birkh\"auser, 1999.
    \item J.-F.\ Le Gall and Y.\ Le Jan, \textit{Branching processes in L\'evy processes: the exploration process}, Ann.\ Probab.\ 26 (1998), 213--252.
\end{enumerate}

\subsection*{Expository Works}
\begin{enumerate}
    \item J.-F.\ Le Gall, \textit{Random geometry on the sphere}, Proc.\ ICM 2014.
    \item J.-F.\ Le Gall, \textit{Brownian motion, random trees, and random surfaces}, lecture notes, 2011.
    \item G.\ Miermont, \textit{Aspects of random maps}, lecture notes, 2014.
\end{enumerate}

